\begin{lemma}[Outside-chain stability at the hierarchy base]
Fix a finite loopless multigraph $G$, a palette $[k]$, a distinguished
edge $e_0$, and an input partial $k$-edge-coloring.  Let
$(T,\varphi_n)$ carry an \emph{ExtendedTashkinovTree} certificate from
that input coloring, and write $T_n$ and $D_n$ for the terminal ladder
tree and its stored-color set.  Let
$\mathcal B$ be a \emph{HierarchyBase} for this certificate under
$\varphi_n$, write
\[
 T_{n,0}^{*}=\mathcal B.\mathit{core}
 \quad\hbox{and}\quad
 R_n=\mathcal B.\mathit{closure},
\]
and suppose that $T$ carries a \emph{TashkinovHierarchy} based at
$\mathcal B$.  Assume that $\varphi_n$ is proper and leaves exactly
$e_0$ uncolored.  If $\sigma_n$ is obtained from $\varphi_n$ by an
outside $(\alpha,\beta)$ Kempe recoloring relative to $T$, then
$\sigma_n$ is proper, leaves exactly $e_0$ uncolored, is
$(T,D_n,\varphi_n)$-stable, and is $(T,\varphi_n)$-invariant.
Furthermore, if $\Theta_n=PE$, then $\sigma_n$ is
$(T_n\oplus R_n,D_n,\varphi_n)$-stable.
\end{lemma}
\begin{proof}
The hypotheses of \noderef{OutsideKempeRecoloringStable} are available
verbatim: $\varphi_n$ is proper, it leaves exactly $e_0$ uncolored, and
the asserted change from $\varphi_n$ to $\sigma_n$ is an outside
$(\alpha,\beta)$ recoloring relative to $T$.  Apply that lemma with
control set $C=D_n$.  Its first four conclusions say respectively that
$\sigma_n$ is proper, leaves exactly $e_0$ uncolored, is
$(T,D_n,\varphi_n)$-stable, and is $(T,\varphi_n)$-invariant.  Retain
also its last conclusion
\[
 \sigma_n(f)=\varphi_n(f)
 \quad\text{for every edge identity $f$ incident with $T$.}       \tag{1}
\]
This proves all conclusions before the final implication.

Assume henceforth that $\Theta_n=PE$.  The parallel clause in
\noderef{HierarchyBase} supplies, in particular, a certificate that
$T_{n,0}^{*}$ is the ordered tree union of $T_n$ and $R_n$.  By
\noderef{OrderedTreeUnion}, this certificate gives
\[
 T_n\text{ is an ordered initial segment of }T_{n,0}^{*},
 \qquad
 V(T_{n,0}^{*})=V(T_n)\cup V(R_n).                    \tag{2}
\]
The core-inside-first clause of \noderef{TashkinovHierarchy} says that
$T_{n,0}^{*}$ is an ordered initial segment of the first hierarchy
stage.  Every subsequent stage is an ordered initial segment of the
next one, and the terminal-stage clause identifies the last stage with
$T$.

We spell out the consequence of this chain.  By
\noderef{TreeSequenceInitialSegment}, each one-step relation consists of
a length inequality together with equality of the vertex and edge entry
at every carried index.  Composing two consecutive steps preserves the
length inequality by transitivity of the numerical order and preserves
each carried entry by transitivity of equality.  Induction over the
finite hierarchy-stage chain therefore shows that
$T_{n,0}^{*}$ is an ordered initial segment of $T$.  In particular, if
$x\in V(T_{n,0}^{*})$, then the description of vertex sets in
\noderef{TreeSequenceVertexSet} writes $x$ as a vertex entry of
$T_{n,0}^{*}$; the initial-segment equality at the same index writes it
as a vertex entry of $T$.  Consequently
\[
 V(T_{n,0}^{*})\subseteq V(T).                         \tag{3}
\]
Combining (2) and (3) gives
\[
 V(T_n)\subseteq V(T),\qquad V(R_n)\subseteq V(T).      \tag{4}
\]

Now let $f$ be incident with $T_n$.  Unfolding
\noderef{TreeIncidentEdge}, there is a vertex $v\in V(T_n)$ incident
with $f$.  The first containment in (4) puts $v$ in $V(T)$, so $f$ is
incident with $T$.  The same argument using the second containment
shows that every edge incident with $R_n$ is incident with $T$.  Thus
(1) restricts to every edge incident with either $T_n$ or $R_n$.

It remains to check the two stable-coloring assertions.  For $T_n$ and
control set $D_n$, the edge clause of \noderef{StableColoring} follows
from (1), even without using its color-side condition.  If
$v\in V(T_n)$, then (4) puts $v$ in $V(T)$, and the vertex clause of the
already obtained $(T,D_n,\varphi_n)$-stability gives
$\overline{\sigma_n}(v)=\overline{\varphi_n}(v)$.  Hence $\sigma_n$ is
$(T_n,D_n,\varphi_n)$-stable.  The identical argument for $R_n$, with
empty control set, proves that $\sigma_n$ is
$(R_n,\varnothing,\varphi_n)$-stable.  These two assertions are exactly
the two conjuncts in \noderef{CombinedStableColoring}; therefore
$\sigma_n$ is $(T_n\oplus R_n,D_n,\varphi_n)$-stable, as required.
\end{proof}
